% ============================================
% ТЕРМИНЫ И ОПРЕДЕЛЕНИЯ
% Оформлено по ГОСТ 7.32 п.4.9: столбцом без знаков препинания в конце
% строк; слева без абзацного отступа в алфавитном порядке — термины,
% справа через тире — определения.
% ============================================

\begingroup
\setlength{\parindent}{0pt}
\setlength{\parskip}{0pt}
\newcommand{\term}[2]{%
	\noindent #1~-- #2\par
	\vspace{0.4em}
}

\term{Глубокий сон}{режим энергопотребления микроконтроллера, при котором основные вычислительные ядра, цифровая периферия и беспроводные модули обесточены, а активными остаются домен реального времени и сопроцессор сверхмалой мощности; пробуждение инициируется аппаратным событием на выводе RTC GPIO, таймером RTC или сигналом сопроцессора ULP}

\term{Декомпозиция энергобюджета}{разложение интегральной энергии цикла обработки по вкладам отдельных этапов, позволяющее идентифицировать доминирующие компоненты и определить приоритеты дальнейшей оптимизации}

\term{Детектор голосовой активности}{компонент аудиосистемы, отличающий периоды наличия речевого сигнала от периодов тишины или фонового шума при минимальных вычислительных затратах; в каскадной архитектуре служит первой ступенью~--- фильтром-предусловием для активации последующих ступеней}

\term{Двухдоменный измерительный стенд}{архитектура лабораторной установки, в которой измеряемая система и система измерений электрически изолированы и питаются от независимых источников, что обеспечивает отсутствие паразитного влияния системы измерений на регистрируемое потребление}

\term{Каскадная архитектура детекции}{организация системы непрерывного мониторинга в виде последовательности ступеней с возрастающей вычислительной сложностью и энергопотреблением, при которой каждая ступень активирует следующую только при срабатывании по своему критерию}

\term{Квантизация нейронной сети}{процедура снижения разрядности представления весов и активаций модели, как правило~--- переход от 32-битного представления с плавающей точкой к 8-битному целочисленному}

\term{Мел-частотные кепстральные коэффициенты}{спектрально-временно\'{е} представление аудиосигнала, получаемое последовательным применением кратковременного преобразования Фурье, проекции амплитудного спектра на психоакустически мотивированную мел-шкалу, логарифмирования и дискретного косинусного преобразования}

\term{Распознавание ключевых слов}{задача обнаружения заранее определённых слов-команд в непрерывном аудиопотоке без полной транскрипции речи; частный случай распознавания речи с ограниченным словарём}

\term{Свёрточная сеть с разделимыми по глубине свёртками}{архитектура свёрточной нейронной сети, в которой стандартная двумерная свёртка по входному тензору заменена двухступенчатой декомпозицией~--- пространственной свёрткой по каждому каналу независимо и линейной комбинацией каналов с ядром~$1 \times 1$~--- что обеспечивает сокращение числа параметров и операций умножения с накоплением в 8--10~раз по сравнению со стандартной свёрткой при сопоставимой точности классификации}

\term{Тензорная арена}{статически выделяемая область оперативной памяти микроконтроллера, используемая фреймворком TensorFlow Lite Micro для размещения промежуточных тензоров нейросетевого инференса без обращения к динамической аллокации}

\endgroup
